\section{The natural signed dual}
\label{sec:signed-dual}

For \(A\Subset\Lambda\), \(\Lcal\chi_A\) is a signed linear combination of
monomials. Writing each coefficient as its sign times its magnitude produces a
finite-set jump process; the remaining diagonal term is represented by a
Feynman--Kac potential.

\subsection{The signed additive set process}
\label{subsec:signed-set-process}

A signed active set is a pair
\[
Y=(A,\sigma)
\in
\{A:A\Subset\Lambda\}\times\{+,-\}.
\]
Signs are multiplied in the usual way and act on real numbers by
\(+x=x\) and \(-x=-x\).

Fix rates and sign labels
\[
\delta_i(S),\ \beta_i(S)\geq0,
\qquad
\sigma_i^\delta(S),\ \sigma_i^\beta(S)\in\{+,-\},
\qquad
 i\in\Lambda,\quad S\subseteq N(i).
\]
We use the convention \(\beta_i(\varnothing)=0\). For \(i\in A\), define
\begin{align}
D_{i,S}(A,\sigma)
&=
\bigl((A\setminus\{i\})\cup S,
\sigma\sigma_i^\delta(S)\bigr),
\label{eq:dual-death-split-map}
\\
B_{i,S}(A,\sigma)
&=
\bigl(A\cup S,
\sigma\sigma_i^\beta(S)\bigr).
\label{eq:dual-birth-map}
\end{align}
The update \(D_{i,\varnothing}\) is called a \emph{death};
\(D_{i,S}\) with \(S\neq\varnothing\) is a \emph{split}; and
\(B_{i,S}\) with \(S\neq\varnothing\) is a \emph{birth}.

The signed additive set process has generator
\begin{equation}
\label{eq:signed-set-generator}
\begin{aligned}
\Dcal g(A,\sigma)
={}&
\sum_{i\in A}\sum_{S\subseteq N(i)}
\delta_i(S)
\bigl(g(D_{i,S}(A,\sigma))-g(A,\sigma)\bigr)
\\
&+
\sum_{i\in A}
\sum_{\substack{S\subseteq N(i)\\S\neq\varnothing}}
\beta_i(S)
\bigl(g(B_{i,S}(A,\sigma))-g(A,\sigma)\bigr).
\end{aligned}
\end{equation}

\subsection{Dual data from the spin rates}
\label{subsec:rates-to-dual}

For \(r\in\R\), let \(\operatorname{sgn}_{\pm}(r)\) be \(+\) when
\(r\geq0\) and \(-\) when \(r<0\). For the coefficients in
\eqref{eq:rate-expansion}, set
\begin{equation}
\label{eq:signed-dual-coefficients}
a_i^\delta(S)=c_i^0(S),
\qquad
a_i^\beta(S)=-c_i^0(S)-c_i^1(S).
\end{equation}
The death and split data are
\begin{equation}
\label{eq:death-split-data}
\delta_i(S)=\abs{a_i^\delta(S)},
\qquad
\sigma_i^\delta(S)=\operatorname{sgn}_{\pm}(a_i^\delta(S)).
\end{equation}
For \(S\neq\varnothing\), the birth data are
\begin{equation}
\label{eq:birth-data}
\beta_i(S)=\abs{a_i^\beta(S)},
\qquad
\sigma_i^\beta(S)=\operatorname{sgn}_{\pm}(a_i^\beta(S)),
\end{equation}
and we set \(\beta_i(\varnothing)=0\) and
\(\sigma_i^\beta(\varnothing)=+\).

The coefficient \(a_i^\beta(\varnothing)\) corresponds to an empty-target
birth, which leaves the active set unchanged. It is therefore included in the
potential rather than represented as a jump. Define
\begin{equation}
\label{eq:dual-potential}
V(A)=\sum_{i\in A}V_i,
\qquad
V_i=
\sum_{S\subseteq N(i)}\delta_i(S)
+
\sum_{\substack{S\subseteq N(i)\\S\neq\varnothing}}\beta_i(S)
+
a_i^\beta(\varnothing).
\end{equation}
The bounds in Section~\ref{sec:spin-systems} imply uniform bounds on the total
jump rate per active site, on \(\abs{V_i}\), and on the number of sites created
by one update.

\subsection{Monomial Feynman--Kac duality}
\label{subsec:monomial-duality}

For \(Y=(A,\sigma)\), define
\begin{equation}
\label{eq:monomial-duality-function}
H(Y,\eta)=\sigma\chi_A(\eta).
\end{equation}
We write \(\pP_A\) and \(\E_A\) for the law and expectation of the signed
set process started from \((A,+)\).

\begin{theorem}[Monomial Feynman--Kac duality]
\label{thm:monomial-duality}
For every signed active set \(Y=(A,\sigma)\),
\begin{equation}
\label{eq:monomial-generator-duality}
\Lcal_\eta H(Y,\eta)
=
\Dcal H(Y,\eta)+V(A)H(Y,\eta).
\end{equation}
Consequently, for every \(A\Subset\Lambda\), \(t\geq0\), and
\(\eta\in\Omega\),
\begin{equation}
\label{eq:monomial-feynman-kac-duality}
P_t\chi_A(\eta)
=
\E_A\left[
\sigma_t
\exp\left(\int_0^t V(A_s)\,ds\right)
\chi_{A_t}(\eta)
\right].
\end{equation}
\end{theorem}

The coefficient calculation proving \eqref{eq:monomial-generator-duality} is
given in Appendix~\ref{app:dual-calculation}. The semigroup identity is the
standard Feynman--Kac consequence of the generator identity.

Independent transitions \(1\to0\) act diagonally on monomials. Let
\begin{equation}
\label{eq:pure-death-noise}
\Ncal^0 f(\eta)
=
\sum_{i\in\Lambda}d_i
\bigl(f(\eta^{i,0})-f(\eta)\bigr),
\qquad d_i\geq0.
\end{equation}
Then
\begin{equation}
\label{eq:pure-death-eigenfunction}
\Ncal^0\chi_A
=
-\left(\sum_{i\in A}d_i\right)\chi_A.
\end{equation}
Thus adding \(\Ncal^0\) leaves the dual jumps unchanged and replaces the
potential \(V(A)\) by
\[
V(A)-\sum_{i\in A}d_i.
\]
We therefore call the independent transitions \(1\to0\) in
\eqref{eq:pure-death-noise} pure deaths.

\subsection{Graphical construction}
\label{subsec:dual-graphical-construction}

For every \(i\in\Lambda\) and \(S\subseteq N(i)\), let
\(I_{i,S}^\delta\subseteq(0,\infty)\) be a Poisson point process of rate
\(\delta_i(S)\). For every nonempty \(S\subseteq N(i)\), let
\(I_{i,S}^\beta\subseteq(0,\infty)\) be a Poisson point process of rate
\(\beta_i(S)\). All these processes are independent. The marked Poisson
interaction set is
\begin{equation}
\label{eq:poisson-interaction-set}
\begin{aligned}
I^{\mathrm P}
={}&
\{(i,t,\delta,S):i\in\Lambda,\ S\subseteq N(i),\ t\in I_{i,S}^\delta\}
\\
&\cup
\{(i,t,\beta,S):i\in\Lambda,\ \varnothing\neq S\subseteq N(i),\ t\in I_{i,S}^\beta\}.
\end{aligned}
\end{equation}
For \((i,t,\alpha,S)\in I^{\mathrm P}\), the site \(i\) is the
\emph{source}, \(S\) is the \emph{target}, and
\(\alpha\in\{\delta,\beta\}\) is the interaction kind. Deaths have empty
target; splits and births have nonempty target.

For an initial state \(Y_0=(A_0,\sigma_0)\), adjoin the deterministic initial
interaction
\begin{equation}
\label{eq:initial-interaction}
\iota_0
=
(\infty,0,\mathsf{init},A_0;\sigma_0),
\end{equation}
where \(\infty\notin\Lambda\) is a formal source. Its target is \(A_0\), its
sign is \(\sigma_0\), and it has no rate. The full interaction set is
\begin{equation}
\label{eq:full-interaction-set}
I=\{\iota_0\}\cup I^{\mathrm P}.
\end{equation}

Start from \(Y_{0-}=(\varnothing,+)\) and apply \(\iota_0\), giving
\(Y_0=(A_0,\sigma_0)\). Thereafter read the interactions in
\(I^{\mathrm P}\) in increasing time order. At a mark
\((i,t,\delta,S)\), apply \(D_{i,S}\) when \(i\in A_{t-}\), and otherwise
do nothing. At a mark \((i,t,\beta,S)\), apply \(B_{i,S}\) when
\(i\in A_{t-}\), and otherwise do nothing. Between interactions the process
is constant.

The uniform bounds noted after \eqref{eq:dual-potential} imply local finiteness:
from a finite initial active set, only finitely many relevant interactions occur
on every bounded time interval. We use the interaction set \(I\) throughout the
patch construction.
